% author: Fabian Bosshard % © CC BY 4.0
\documentclass[9pt,headings=standardclasses,parskip=half]{scrartcl}

% ===================== basics =====================
\usepackage{ifthen,iftex,csquotes}
\usepackage[T1]{fontenc}
\usepackage[english]{babel}

% ===================== graphics / color =====================
\usepackage{graphicx}
\usepackage[dvipsnames]{xcolor}
\usepackage[left=38mm,right=38mm,top=20mm,bottom=30mm]{geometry}

% ===================== math =====================
\usepackage{amsmath,amssymb,amsthm,mathtools,mathdots,accents}

% ===================== symbols =====================
\usepackage{pifont}
\usepackage{marvosym}

% ===================== TikZ / plots =====================
\usepackage{tikz}
\usetikzlibrary{arrows,arrows.meta,shapes,positioning,calc,fit,patterns,intersections,math,3d,tikzmark,decorations.pathreplacing,decorations.markings}
\usepackage{tikz-dependency,tikz-qtree,tikz-qtree-compat,tikz-3dplot,tikzpagenodes}
\usepackage{pgfplots}




\makeatletter
% knob: multiply font rule thickness by MUL/DIV (integers only)
\newcount\matrixsymbolthickmul \matrixsymbolthickmul=10   % numerator
\newcount\matrixsymbolthickdiv \matrixsymbolthickdiv=10   % denominator
\newcount\matrixsymboloffsetdiv \matrixsymboloffsetdiv=10 % vertical offset = ht(box)/this

% \matrixsymbol@ takes 5 args:
%   #1 = math style token (\displaystyle, \textstyle, \scriptstyle, \scriptscriptstyle)
%   #2 = a math font (e.g. \textfont3, \scriptfont3, \scriptscriptfont3) used to read \fontdimen8
%   #3 = the actual symbol/formula to typeset (usually #1 from \matrixsymbol)
%   #4 = trim (each side) to subtract from underline width (left & right)
%   #5 = extend (total) to add to underline width (centered extension)
\protected\def\matrixsymbol@#1#2#3#4#5{% 
    % Typeset the symbol in the chosen math style and store it in box0
    \setbox0=\hbox{$\m@th#1#3$}
    %
    % Rule thickness: take the font's default rule thickness (fontdimen8) for the chosen font #2, then scale it by thickmul/thickdiv
    \dimen6=\fontdimen8#2\relax
    \multiply\dimen6 by \matrixsymbolthickmul
    \divide\dimen6 by \matrixsymbolthickdiv
    %
    % Vertical placement: fixed below the baseline as a fraction of the symbol box height
    % (distance = ht(box0)/matrixsymboloffsetdiv)
    \dimen8=\ht0 \divide\dimen8 by \matrixsymboloffsetdiv
    %
    % Underline width computation:
    % start with symbol width
    \dimen0=\wd0 
    % add total extension (#5)
    \advance\dimen0 by #5\relax
    % subtract 2 * trim (#4) because trim is per-side
    \dimen2=#4\relax 
    \advance\dimen0 by -2\dimen2
    %
    % Horizontal start shift for the underline relative to the symbol:
    % We want the extension (#5) split equally left/right (so shift left by #5/2), then apply trim (#4) (shift right by trim).
    % net shift = trim - extend/2
    \dimen4=#5\relax \divide\dimen4 by 2 \dimen4=-\dimen4 \advance\dimen4 by \dimen2
    %
    % Output: a box whose visible content is copy0, plus an overlaid underline:
    % - \rlap makes the underline overlay without consuming horizontal space
    % - \smash makes it add no height/depth (so it won't affect vertical spacing)
    % - \raise moves it below the baseline by dimen8
    \hbox{\rlap{\smash{\raise-\dimen8\hbox{\kern\dimen4\vrule width\dimen0 height0pt depth\dimen6}}}\copy0}%
}
\DeclareRobustCommand{\matrixsymbol}[1]{%
    \begingroup
    \let\@nomath\@gobble \mathversion{bold}% make #1 bold by switching to the bold math version
    %
    % \math@atom{#1}{<stuff>} preserves the *atom type* of #1 (ord/op/bin/rel/...)
    % so spacing and script placement behave like the original symbol.
    \math@atom{#1}{%
    \mathchoice
    {\matrixsymbol@{\displaystyle}{\textfont3}{#1}{0.03em}{0.00em}}% 
    {\matrixsymbol@{\textstyle}{\textfont3}{#1}{0.03em}{0.00em}}%
    {\matrixsymbol@{\scriptstyle}{\scriptfont3}{#1}{0.03em}{0.00em}}%
    {\matrixsymbol@{\scriptscriptstyle}{\scriptscriptfont3}{#1}{0.03em}{0.00em}}%
    }%
    \endgroup
}
\makeatother

\newcommand{\matr}[1]{\matrixsymbol{#1}}

\begin{document}
\pagestyle{empty}


\[
  \matr{A} = \matr{Q}\matr{R} \boldsymbol{x}
\]
\[
  \underline{\boldsymbol{A}} = \underline{\boldsymbol{Q}}\underline{\boldsymbol{R}} \boldsymbol{x}
\]
\[
  \boldsymbol{A} = \boldsymbol{Q}\boldsymbol{R} \boldsymbol{x}
\]
\[
  e^{\matr{A}} \neq e^{\matr{Q}}e^{\matr{R}}
\]
\[
  e^{\underline{\boldsymbol{A}}} \neq e^{\underline{\boldsymbol{Q}}}e^{\underline{\boldsymbol{R}}}
\]
\[
  e^{\boldsymbol{A}} \neq e^{\boldsymbol{Q}}e^{\boldsymbol{R}}
\]
\[
  \matr{A}_1  \matr{Q}_1 \matr{R}_1
\]
\[
  \boldsymbol{A}_1 \boldsymbol{Q}_1 \boldsymbol{R}_1
\]
\[
\matr{A}\matr{Q}\matr{W}\matr{R}\matr{I}
\]
\[
  e^{\matr{A}\matr{Q}\matr{W}\matr{R}\matr{I}}
\]
\[
  e^{e^{\matr{A}\matr{Q}\matr{W}\matr{R}\matr{I}}}
\]
\[
  e^{e^{\boldsymbol{A}\boldsymbol{Q}\boldsymbol{W}\boldsymbol{R}\boldsymbol{I}}}
\]
\[
  \matr{A} = \matr{W}\matr{Q}
\]
\[
  \matr{A} = \matr{Q}\matr{R}
\]
\[
  \matr{A} = \matr{I}\matr{I}
\]
\[
\tilde{\matr{\chi}}_1^{(1)} \tilde{\chi}_1^{(1)} \tilde{b}_1^{(1)}
\]
\[
\matr{\chi}_1^{(1)} \chi_1^{(1)} b_1^{(1)}
\]




% // it seems to work...

% // but it doesnt...


% #let matr1(x) = context {
%   let em_abs = measure(math.inline[box(width: 1em)]).width.to-absolute()
%   let mu = em_abs / 18
%   let matrspace = 0.1 * mu

%   math.class("normal",
%     h(matrspace)
%     + math.underline(
%         math.squish(h(-matrspace) + math.bold(x) + h(-matrspace), mode: "bottom")
%       )
%     + h(matrspace)
%   )
% }

% #let matr2(x) = context {
%   let em_abs = measure(math.inline[box(width: 1em)]).width.to-absolute()
%   let mu = em_abs / 18
%   let matrspace = 0.1 * mu

%   math.class("normal", math.squish(
%     h(matrspace)
%     + math.underline(
%         math.squish(h(-matrspace) + math.bold(x) + h(-matrspace), mode: "bottom")
%       )
%     + h(matrspace), mode: "bottom")
%   )
% }



% // compare dfferences between the approaches
% $ 
% matr1(A)_(matr1(A)_matr1(A))^matr1(A)^matr1(A) // subscripts do not work correctly
% matr2(A)_(matr2(A)_matr2(A))^matr2(A)^matr2(A) // subscripts work correctly
% bold(A)_(bold(A)_bold(A))^bold(A)^bold(A) // subscripts work correctly
% bold(underline(A))_(bold(underline(A))_bold(underline(A)))^bold(underline(A))^bold(underline(A)) // subscripts work correctly (in terms of underline philosphy)
% $


% // compare vertical spacing correctness
% $
% matr2(A)_(matr2(A)_matr2(A))^matr2(A)^matr2(A) // subscripts work correctly
% bold(A)_(bold(A)_bold(A))^bold(A)^bold(A) // subscripts work correctly
% $



% // compare horizontal spacing correctness
% $
% matr2(A)_(matr2(A)_matr2(A))^matr2(A)^matr2(A) // subscripts work correctly
% $
% $
% bold(A)_(bold(A)_bold(A))^bold(A)^bold(A) // subscripts work correctly
% $






% // compare dfferences between the approaches
% $ 
% matr1(Q)_(matr1(Q)_matr1(Q))^matr1(Q)^matr1(Q) // subscripts do not work correctly
% matr2(Q)_(matr2(Q)_matr2(Q))^matr2(Q)^matr2(Q) // subscripts do not work correctly
% bold(Q)_(bold(Q)_bold(Q))^bold(Q)^bold(Q) // subscripts work correctly
% bold(underline(Q))_(bold(underline(Q))_bold(underline(Q)))^bold(underline(Q))^bold(underline(Q)) // subscripts work correctly (in terms of underline philosphy)
% $


% // compare vertical spacing correctness
% $
% matr2(Q)_(matr2(Q)_matr2(Q))^matr2(Q)^matr2(Q) // subscripts do not work correctly
% bold(Q)_(bold(Q)_bold(Q))^bold(Q)^bold(Q) // subscripts work correctly
% $



% // compare horizontal spacing correctness
% $
% matr2(Q)_(matr2(Q)_matr2(Q))^matr2(Q)^matr2(Q) // subscripts do not work correctly
% $
% $
% bold(Q)_(bold(Q)_bold(Q))^bold(Q)^bold(Q) // subscripts work correctly
% $


% compare vertical spacing
\[
\matr{A}_{\matr{A}_{\matr{A}}}^{\matr{A}^{\matr{A}}}
\boldsymbol{A}_{\boldsymbol{A}_{\boldsymbol{A}}}^{\boldsymbol{A}^{\boldsymbol{A}}}
\]
% compare  horizontal spacing
\[
\matr{A}_{\matr{A}_{\matr{A}}}^{\matr{A}^{\matr{A}}}
\]
\[
\boldsymbol{A}_{\boldsymbol{A}_{\boldsymbol{A}}}^{\boldsymbol{A}^{\boldsymbol{A}}}
\]

%now with Q
% compare vertical spacing
\[
\matr{Q}_{\matr{Q}_{\matr{Q}}}^{\matr{Q}^{\matr{Q}}}
\boldsymbol{Q}_{\boldsymbol{Q}_{\boldsymbol{Q}}}^{\boldsymbol{Q}^{\boldsymbol{Q}}}
\]
% compare  horizontal spacing
\[
\matr{Q}_{\matr{Q}_{\matr{Q}}}^{\matr{Q}^{\matr{Q}}}
\]
\[
\boldsymbol{Q}_{\boldsymbol{Q}_{\boldsymbol{Q}}}^{\boldsymbol{Q}^{\boldsymbol{Q}}}
\] 



\end{document}